\documentclass{beamer}

\usetheme{Madrid}
\usecolortheme{default}

\usepackage[italian]{babel}
\usepackage[utf8]{inputenc}
\usepackage[T1]{fontenc}
\usepackage{xcolor}
\usepackage{colortbl}
\usepackage{verbatim}

\title{Basi di Dati}
\subtitle{SQL DML -- Manipolazione dei dati}
\author{Prof. Federico Mazzini, Prof.ssa Francesca Larotonda}
\institute{I.I.S. Francesco Alberghetti}
\date{A.S. 2025 -- 2026}

\begin{document}

%==============================================================================
\begin{frame}
\titlepage
\end{frame}

%==============================================================================
\section{SQL DML (Data Manipulation Language)}
%==============================================================================

\begin{frame}{Dal DDL al DML}
Con il DDL abbiamo definito la \textbf{struttura} del database:
\begin{itemize}
    \item tabelle;
    \item colonne;
    \item tipi di dato;
    \item vincoli di integrità.
\end{itemize}

\medskip
Per lavorare sui \textbf{record} del database invece, utilizziamo il DML (\textit{Data Manipulation Language})
\begin{itemize}
    \item inserimento
    \item modifica
    \item cancellazione
\end{itemize}
\end{frame}

%==============================================================================
\begin{frame}{Cos'è il DML}
\begin{block}{Definizione}
Il DML (Data Manipulation Language) è la parte di SQL che permette di \textbf{inserire, modificare ed eliminare dati} all'interno delle tabelle.
\end{block}

\medskip
Principali comandi DML:
\begin{itemize}
    \item \texttt{INSERT}
    \item \texttt{UPDATE}
    \item \texttt{DELETE}
\end{itemize}

\medskip
Il DML agisce sui \textbf{record}, non sulla struttura.
\end{frame}

%==============================================================================
\section{INSERT}
%==============================================================================


%==============================================================================
\begin{frame}[fragile]{INSERT}

Il comando \textbf{INSERT} serve per inserire nuovi record in una tabella.
\begin{block}{Sintassi}
\begin{verbatim}
INSERT INTO nome_tabella (colonna1, colonna2, ...)
VALUES (valore1, valore2, ...);
\end{verbatim}
\end{block}

\end{frame}

%==============================================================================
\begin{frame}[fragile]{INSERT}
\small

\textbf{Inserimento di uno studente}

\medskip
\texttt{studente(\underline{id\_studente}, matricola, nome, cognome, data\_nascita)}

\begin{verbatim}
INSERT INTO studente (matricola, nome, cognome, data_nascita)
VALUES (1001, 'Luca', 'Rossi', '2007-03-15');
\end{verbatim}

\medskip
\textbf{Inserimento di un professore}

\medskip
\texttt{professore(\underline{id\_professore}, cod\_fiscale, nome, cognome)}

\begin{verbatim}
INSERT INTO professore (cod_fiscale, nome, cognome)
VALUES ('BNCMRA80A01H501U', 'Marco', 'Bianchi');
\end{verbatim}

\medskip
Le colonne id\_* sono AUTO\_INCREMENT, quindi non vanno specificate
nell'INSERT: il DB le genera automaticamente.
\end{frame}

%==============================================================================
\begin{frame}[fragile]{INSERT singola VS. multipla}
\small

\medskip
\textbf{Due INSERT separate:}

\begin{verbatim}
INSERT INTO studente (matricola, nome, cognome, data_nascita)
VALUES (1002, 'Anna', 'Bianchi', '2007-07-21');

INSERT INTO studente (matricola, nome, cognome, data_nascita)
VALUES (1003, 'Paolo', 'Verdi', '2006-11-03');
\end{verbatim}

\medskip
\textbf{Un'unica INSERT con più record:}

\begin{verbatim}
INSERT INTO studente (matricola, nome, cognome, data_nascita)
VALUES (1004, 'Giulia', 'Neri', '2007-02-10'),
       (1005, 'Marco', 'Galli', '2006-09-18');
\end{verbatim}

\end{frame}

%==============================================================================
\section{UPDATE}


%==============================================================================
\begin{frame}[fragile]{UPDATE}

Il comando \textbf{UPDATE} modifica uno o più record esistenti.
\begin{block}{Sintassi}
\begin{verbatim}
UPDATE nome_tabella
SET colonna1 = valore1,
    colonna2 = valore2
WHERE condizione;
\end{verbatim}
\end{block}
\end{frame}

%==============================================================================
\begin{frame}[fragile]{UPDATE}
\small
\textbf{Modifica del cognome di uno studente:}

\begin{verbatim}
UPDATE studente
SET cognome = 'Bianchi'
WHERE id_studente = 1;
\end{verbatim}

\medskip
La clausola WHERE identifica quali record aggiornare.

\medskip
\alert{Cosa succede se la clausola \texttt{WHERE} non viene specificata?}
\end{frame}

%==============================================================================
\begin{frame}[fragile]{UPDATE}
\small

\begin{block}{Clausola Where}
La clausola \textbf{WHERE} non identifica necessariamente
un solo record, ma \textbf{tutti quelli che soddisfano la condizione}.
\end{block}

\medskip
\textbf{Esempio: modificare tutti gli studenti che si chiamano Luca}

\begin{verbatim}
UPDATE studente
SET cognome = 'Verdi'
WHERE nome = 'Luca';
\end{verbatim}

\medskip
Se nel database esistono 3 studenti con nome = 'Luca',
tutti e 3 verranno aggiornati.

\end{frame}

%==============================================================================
\section{DELETE}
%==============================================================================
\begin{frame}[fragile]{DELETE}

Il comando \textbf{DELETE} elimina uno o più record da una tabella.
\begin{block}{Sintassi}
\begin{verbatim}
DELETE FROM nome_tabella
WHERE condizione;
\end{verbatim}
\end{block}
\end{frame}

%==============================================================================
\begin{frame}[fragile]{DELETE}
\small
\begin{verbatim}
DELETE FROM studente
WHERE id_studente = 1;
\end{verbatim}

\medskip
Il record con id\_studente 1 viene eliminato dalla tabella.
\end{frame}

%==============================================================================
\begin{frame}[fragile]{DELETE senza WHERE}
\small
\begin{verbatim}
DELETE FROM studente;
\end{verbatim}

\medskip
\alert{Elimina TUTTI gli studenti}

\end{frame}

%==============================================================================
\begin{frame}[fragile]{Vincoli da rispettare nelle DML}
\small

Ogni operazione DML (\texttt{INSERT}, \texttt{UPDATE}, \texttt{DELETE}) deve rispettare i vincoli strutturali delle tabelle. Il database mantiene sempre l'integrità dei dati.

\medskip
\textbf{Vincoli principali:}

\begin{itemize}
    \item \texttt{NOT NULL} → non è possibile inserire valori nulli.
    \item \texttt{PRIMARY KEY e UNIQUE} → non è possibile inserire valori duplicati.
    \item \texttt{FOREIGN KEY} → non è possibile eliminare record referenziati, o assegnare record a elementi inesistenti.
\end{itemize}

\medskip
\alert{Il DBMS interviene automaticamente per proteggere la coerenza dei dati.}

\end{frame}

%==============================================================================

\end{document}
